\documentclass[12pt]{article}
\usepackage[utf8]{inputenc}
\usepackage[T1]{fontenc}
\usepackage{hyperref}
\title{Kahayan Language Resources for K-12 Curriculum Development}
\author{Compiled by Claude Sonnet 5.5 (claude-sonnet-5-5) for Andr\'as Kornai}
\date{}
\begin{document}
\maketitle

\section{Introduction}
Kahayan (ISO 639-3 xah, Glottocode kaha1240) is spoken in Central Kalimantan, Indonesia. The identification line in the request gives the family as ``Bookkeeping''. That is evidently a data error and was not used. The search results place Kahayan in the Austronesian language area, as a dialect or close relative of Ngaju. They disagree on its status: one source lists Kahayan among Ngaju dialects, another names only Pulau Petak, Ba'amang and Mantangai. Almost all resources found are for Ngaju in general, not Kahayan specifically. No Kahayan-specific school materials turned up in the searches.

\section{Primary Educational Resources}
A web-based Ngaju--English picture dictionary application for primary-school children is described in a conference paper from UIN Palangka Raya:
\url{https://e-proceedings.uin-palangkaraya.ac.id/INACELT/article/download/84/87/145}.
A related journal article is at \url{https://journal.umpr.ac.id/index.php/jsakti/article/view/12580}. I did not open it, so its content is not verified here.

\section{Linguistic and Reference Materials}
\begin{itemize}
\item Overview of the Ngaju language and its dialects: \url{https://en.wikipedia.org/wiki/Ngaju_language}
\item Ngaju language profile (Missouri DICE): \url{https://dice.missouri.edu/assets/docs/austronesia/DayakNgaju.pdf}
\item Austronesian Comparative Dictionary entry: \url{https://acd.clld.org/values/18577-4f54e367268784734aec17693bf5790d-1}
\item Open Library record for a Ngaju-related book: \url{https://openlibrary.org/books/OL153128M?v=16}
\item KITLV page: \url{https://kitlv.nl/?p=2174}
\end{itemize}
The search summary also mentioned \emph{Kamus Bahasa Dayak Ngaju -- Indonesia} (Jakarta, 2009, Anugerah Indah Mandiri Press). No URL for it was seen, so it is listed here only and not in the JSON.

\section{Religious and Community Resources}
No Ngaju or Kahayan Bible translation was confirmed in the search results. Martin Baier compiled terms and examples of the traditional priest language (\emph{bahasa sangiang}) with modern Ngaju, Indonesian and German equivalents. No URL for that work was seen.

\section{Digital and Audio Resources}
Verbix lists Ngaju: \url{https://docs.verbix.com/Languages/Ngaju}. No audio or video resources were found.

\section{Curriculum Development Context}
Very little is documented for Kahayan itself. A curriculum would need to start from Ngaju materials, check them with Kahayan speakers, and commission new Kahayan-specific material. No grammar, textbooks or school programmes were found in these searches.
\end{document}